\documentclass[border=10pt,tikz]{standalone}
\usepackage{tikz}
\usepackage{tikz-timing}
\usepackage{fontspec}
\setmainfont{Apple SD Gothic Neo}
\setsansfont{Apple SD Gothic Neo}
\usetikztiminglibrary{counters,interval,nicetabs}
\begin{document}
\begin{tikzpicture}[font=\scriptsize]

\node[font=\footnotesize\bfseries, anchor=west] at (0, 4) {UART Hardware Flow Control — RTS/CTS handshake};

\node[anchor=west] at (0, 2) {\begin{tikztimingtable}[xscale=1.5, yscale=1.4]
  A.RTS  & H 0.5L 3L 0.5H H \\
  B.CTS  & H 0.5L 3L 0.5H H \\
  A.TX   & H 1.5H 1.2D{Hello} 1.2D{World} 1H \\
  A.CTS  & H L 4L L H \\
  B.RTS  & H L 4L L H \\
\end{tikztimingtable}};

\node[anchor=west, font=\tiny, align=left] at (0, -2) {
  \textbullet\ A의 RTS = "보낼 데이터 있음" 신호 (A → B)\\
  \textbullet\ B의 CTS = "받을 준비 됨" 응답 (B → A)\\
  \textbullet\ A·B 양쪽 CTS Low → A는 TX 송신\\
  \textbullet\ B의 버퍼 차오르면 B가 자기 RTS High로 풀어 → A의 CTS High → A 정지\\
  \textbullet\ 핀 — A.RTS↔B.CTS, A.CTS↔B.RTS (cross)
};

\end{tikzpicture}
\end{document}
